% from ega4-14-fr.tex
\oldpage[IV-3]{205}

Or, si $y'=f'(x')$, la projection $f'^{-1}(y')\to f^{-1}(y)$ est un morphisme
ouvert (2.4.10), et il y a, en vertu de (14.1.4), un voisinage $V'$ de $x'$ dans
$f'^{-1}(y')$ où $f'$ n'est pas ouvert~; donc $f$ n'est pas universellement ouvert
aux points de l'image de $V'$ dans $f^{-1}(y)$, qui est un voisinage de $x$ dans
$f^{-1}(y)$.

\begin{proposition}[14.3.4]
\label{IV.14.3.4-fr}
\begin{enumerate}[label=(\roman*)]
  \item Soient $f:X\to Y$, $g:Y\to Z$ deux morphismes, $x$ un point de $X$,
  $y=f(x)$. Si $f$ est universellement ouvert au point $x$ et $g$ universellement
  ouvert au point $y$, alors $g\circ f$ est universellement ouvert au point $x$.
  Inversement, si $g\circ f$ est universellement ouvert au point $x$, $g$ est
  universellement ouvert au point $y$.

  \item Si $f:X\to Y$ est un $S$-morphisme universellement ouvert au point
  $x\in X$, alors, pour tout changement de base $S'\to S$,
  $f_{(S')}:X_{(S')}\to Y_{(S')}$ est universellement ouvert en tout point de
  $X_{(S')}$ au-dessus de $x$.

  \item Pour que $f:X\to Y$ soit universellement ouvert au point $x\in X$, il
  faut et il suffit que $f_{\mathrm{red}}$ le soit.

  \item Soient $f:X\to Y$ un morphisme localement de présentation finie, $x$ un
  point de $X$, $y=f(x)$~; posons $Y_1=\operatorname{Spec}(\mathcal O_y)$,
  $X_1=X\times_Y Y_1$, $f_1=f_{(Y_1)}$. Pour que $f$ soit universellement ouvert
  au point $x$, il faut et il suffit que $f_1$ le soit (on rappelle (I, 3.6.5) que
  $X_1$ est canoniquement identifié à un sous-espace de $X$).
\end{enumerate}
\end{proposition}

\begin{proof}
En effet (ii) est une conséquence évidente de la définition (14.3.3)~; il résulte
aussi de la définition que pour démontrer l'assertion (i), il suffit de le faire
lorsqu'on y supprime partout le mot «~universellement~», et cela résulte alors de
(14.1.2). L'assertion (iii) résulte de (ii) et de ce que le morphisme canonique
$X_{\mathrm{red}}\to X$ est surjectif. Enfin, la condition (iv) est trivialement
nécessaire. D'autre part, si elle est vérifiée, et si $g:Y'\to Y$ est un morphisme
quelconque, $X'=X_{(Y')}$, $f'=f_{(Y')}$, $x'$ un point de $X'$ au-dessus de $x$,
alors $f'$ est localement de présentation finie, et pour voir qu'il est ouvert au
point $x'$, il suffit d'appliquer le critère (1.10.3, $b$). Soient
$y'=f'(x')$, $Y'_1=\operatorname{Spec}(\mathcal O_{y'})$,
$X'_1=X'\times_{Y'}Y'_1$, $f'_1=f_{(Y'_1)}$. Comme $g(y')=y$, le morphisme
composé $Y'_1\to Y'\to Y$ se factorise en $Y'_1\to Y_1\to Y$ (I, 2.4.4), donc
$X'_1=X_1\times_{Y_1}Y'_1$ et $f'_1=(f_1)_{(Y'_1)}$~; la conclusion résulte alors
aussitôt de l'hypothèse que $f'_1$ est ouvert au point $x'$ et de (1.10.3).
\end{proof}

\begin{proposition}[14.3.5]
\label{IV.14.3.5-fr}
Soient $X$, $Y$ deux préschémas, $f:X\to Y$ un morphisme, $x$ un point de $X$.
Soit $(Y_i)_{1\leq i\leq n}$ une famille localement finie de sous-préschémas
fermés de $Y$ telle que l'espace $Y$ soit réunion des $Y_i$, et supposons que pour
tout $i$ tel que $f(x)\in Y_i$, la restriction $f^{-1}(Y_i)\to Y_i$ de $f$ soit un
morphisme universellement ouvert au point $x$~; alors $f$ est universellement
ouvert au point $x$.
\end{proposition}

\begin{proof}
Compte tenu de la définition, cela résulte de la proposition analogue (14.1.2,
(ii)) pour les morphismes ouverts en un point.
\end{proof}

\begin{proposition}[14.3.6]
\label{IV.14.3.6-fr}
Soient $Y$ un préschéma localement noethérien, $f:X\to Y$ un morphisme localement
de type fini. Pour que $f$ soit universellement ouvert en un point $x\in X$, il
faut et il suffit que la condition suivante soit vérifiée~: pour tout morphisme
$g:Y'\to Y$, où $Y'=\operatorname{Spec}(A)$ est le spectre d'un anneau de
valuation discrète, tel que l'image $g(y')$ du point fermé $y'$ de $Y'$ soit égale
à $y=f(x)$, et pour tout point $x'\in X'=X\times_Y Y'$ dont les projections sur
$X$ et $Y'$ sont $x$ et $y'$ respectivement, il existe une générisation $z'$ de
$x'$ dans $X'$ dont la projection dans $Y'$ est le point générique de $Y'$
(autrement dit, il existe une composante irréductible de $X'$ contenant $x'$ et
dominant $Y'$).

